\chapter{Audit TI dan Asurans}

\section*{Tujuan Pembelajaran}
\begin{itemize}
  \item Menjelaskan pembagian peran dalam Three Lines Model dan posisi audit internal yang independen.
  \item Menyusun rencana asurans berbasis risiko untuk program transformasi digital.
  \item Menguji rancangan dan efektivitas satu kontrol dengan bukti yang memadai.
\end{itemize}

\section{Pendahuluan}

% TODO
% Cakupan: direksi menerima laporan "semua aman" dari unit TI, padahal audit menemukan sebaliknya; direksi butuh keyakinan yang independen.

\section{Konsep Asurans}

% TODO
% Cakupan: asurans dan bedanya dengan konsultasi; pihak yang terlibat (pemilik tanggung jawab, penilai, pengguna hasil); asurans internal dan eksternal; tingkat keyakinan yang dapat diberikan.

\section{Three Lines Model}

% TODO
% Cakupan: peran badan pengatur, manajemen (lini pertama dan kedua), dan audit internal (lini ketiga); Statement of Position IIA tanggal 8 Juli 2026 yang menggantikan position paper 2020; independensi dan pengamanannya; kesalahan umum penerapan.
% Kasus Medikara: unit manajemen risiko yang melapor kepada direktur TI.

\section{Global Internal Audit Standards}

% TODO
% Cakupan: lima domain dan 15 prinsip; berlaku sejak 9 Januari 2025; implikasi bagi fungsi audit internal yang kecil; ITAF dari ISACA dan sertifikasi profesi audit TI sebagai rujukan pelengkap.

\section{Perencanaan Audit Berbasis Risiko}

% TODO
% Cakupan: semesta audit; prioritas audit dari risk register dan register kewajiban; rencana asurans untuk program transformasi; koordinasi dengan penyedia asurans lain seperti auditor eksternal dan lembaga sertifikasi.

\section{Program Audit dan Pengujian Kontrol}

% TODO
% Cakupan: tujuan, kriteria, dan prosedur audit; uji rancangan dan uji efektivitas kontrol; teknik sampling; kecukupan dan keandalan bukti.
% Kasus Medikara: peninjauan hak akses SIMRS yang menjadi temuan audit internal.
% Pendalaman teknis: teknik audit berbantuan komputer pada log akses (akun tanpa pemilik, akun bersama, akses di luar jam kerja) dan konsep continuous control monitoring.

\section{Pelaporan Temuan dan Tindak Lanjut}

% TODO
% Cakupan: format kondisi, kriteria, sebab, akibat, dan rekomendasi; peringkat temuan; pemantauan tindak lanjut; memanfaatkan laporan asurans pihak ketiga seperti SOC 2.

\section{Aktivitas: Rencana Asurans dan Pengujian Satu Kontrol}

% TODO
% Tujuan aktivitas: memastikan peta jalan transformasi dapat diperiksa secara independen dan berlatih menguji kontrol dengan bukti.
% Peran dalam tim: anggota berlatar akuntansi atau audit memimpin; analis teknologi menyiapkan dan menganalisis data; pemilik konteks menilai apakah rekomendasi dapat dijalankan; analis risiko menautkan rencana asurans ke risk register.
% Langkah: petakan lini pertama, kedua, dan ketiga di organisasi; susun rencana asurans dua tahun pertama peta jalan; pilih satu kontrol penting; susun program audit singkat; uji kontrol dengan data contoh; tulis temuan dalam format baku.
% Luaran dan tonggak capstone: rencana asurans berbasis risiko dan contoh pengujian kontrol beserta temuan; masuk C4.
% Berkas sources yang direncanakan: sources/chapter12/rencana-asurans.xlsx, sources/chapter12/program-audit-hak-akses.md, sources/chapter12/data-contoh-hak-akses.csv

\section{Ringkasan}

% TODO

\section{Latihan}

\begin{enumerate}
  \item Unit manajemen risiko di sebuah rumah sakit melapor kepada direktur TI. Apa masalahnya menurut Three Lines Model, dan bagaimana memperbaikinya?
  \item Jelaskan perbedaan uji rancangan dan uji efektivitas dengan contoh kontrol peninjauan hak akses.
  \item Tulis satu temuan audit lengkap untuk kasus kebocoran data lewat vendor Medikara pada 2025.
\end{enumerate}
